\subsection{Decisions from solvable finite congruence problems}
\label{C:Greenberg-congruence-subsection}

\begin{lemma}[Proved: integral, projective and factorization decisions]
\label{C:Greenberg-congruence-decisions}
Let $F$ be a countable henselian discretely valued field with
decidable exact finite names, equality and arithmetic, a computable
normalized valuation, and a supplied uniformizer $\pi$.
Put $\mathcal O=\{0\}\cup\{a:v(a)\ge0\}$.
Suppose that, uniformly in $N\ge1$ and a finite list of polynomials
with explicitly supplied coefficients in $\mathcal O$, existence
of a simultaneous zero in $\mathcal O/\pi^N\mathcal O$ is decidable
by an ordinary terminating algorithm.  If $F$ has positive
characteristic, suppose also that its completion $\widehat F/F$
is separable.  No finiteness hypothesis on the residue field is made.

There are ordinary terminating algorithms deciding simultaneous
integral polynomial equations over $\mathcal O$ and rational points
on every explicitly presented projective $F$-scheme.  Every positive
answer constructs an exact solution or point.  There is also ordinary
univariate factorization over $F$, with multiplicities and actual
factor coefficients.  No effective bound for the required congruence
precision is assumed or claimed.
\end{lemma}
\begin{proof}
For polynomials $f_1,\ldots,f_s\in\mathcal O[X_1,\ldots,X_m]$,
write
\[
 Z=\operatorname{Spec}\bigl(\mathcal O[X_1,\ldots,X_m]/(f_1,\ldots,f_s)\bigr).
\]
Greenberg's approximation theorem and its Corollary~2 give
\[
 Z(\mathcal O)\ne\varnothing
 \quad\Longleftrightarrow\quad
 Z(\mathcal O/\pi^N\mathcal O)\ne\varnothing
       \text{ for every }N\ge1 .                    \tag{C.3.GC1}
\]
The hypotheses are exactly a henselian discrete valuation ring
and, in positive characteristic, separability of the completion;
the theorem does not require a finite residue field
\cite[Theorem 1 and Corollary 2, p.\ 59, with proof pp.\ 60--64]
{C-Greenberg1966}.  We use only this existence equivalence, not
computability of the constants in the stronger approximation statement.

Enumerate all exact $F$-tuples, retain the integral tuples by the
valuation test, and test all the equations.  This is a positive
search for an actual point of $Z(\mathcal O)$.
In parallel, for $N=1,2,\ldots$, apply the supplied terminating
congruence decision to the same equations modulo $\pi^N$.
A negative answer at any level is a correct negative answer over
$\mathcal O$.  If no integral solution exists, (C.3.GC1) forces
such a finite level.  If an integral solution exists, the positive
search finds one.  Dovetailing finite computation steps therefore
always terminates with a correct answer and, in the positive case,
an exact tuple.  This also handles empty lists of equations,
constant equations and zero-dimensional affine schemes.

Now let $X\subseteq\mathbf P_F^m$ be given by homogeneous equations
$h_1,\ldots,h_s$.  Multiply each by a power of $\pi$ to make all
its coefficients integral, without changing its $F$-zero locus.
Apply the integral decision just proved to the affine equations
\[
 h_j(x_0,\ldots,x_m)=0\ (1\le j\le s),
 \qquad \sum_{i=0}^m x_i y_i=1.                     \tag{C.3.GC2}
\]
An integral solution has a nonzero $x$-vector and therefore gives
an $F$-point of $X$.  Conversely, scale the coordinates of any
$F$-point by a power of $\pi$ so that their minimum valuation is
zero.  They are then all integral and some $x_i$ is a unit.
Set $y_i=x_i^{-1}$ for one such index and all other $y_j$ equal
to zero.  This is an integral solution of (C.3.GC2).
Thus the affine integral decision is exactly the required projective
decision, including nonreduced and singular schemes.  The same
encoding in each quotient expresses a primitive vector, so no
finite-branching or residue-field compactness assertion is used.

For factorization, first let $f\in\mathcal O[T]$ be monic of degree
$n\ge2$.  For each $1\le d<n$, write monic indeterminates
$g=T^d+\sum_{i<d}a_iT^i$ and
$h=T^{n-d}+\sum_{j<n-d}b_jT^j$.
The coefficient equations for $gh=f$ are a finite polynomial
system over $\mathcal O$.  Apply the integral decision to each.
Every monic factor of $f$ over $F$ lies in $\mathcal O[T]$:
its roots are integral over $\mathcal O$, hence so are its
coefficients, and $\mathcal O$ is integrally closed in $F$.
Consequently all answers are negative exactly when $f$ is
irreducible in $F[T]$.  A positive answer constructs genuine monic
factors of strictly smaller degree.  Recursion and exact division
give the full factorization and all multiplicities, including
inseparable factors.

For an arbitrary nonconstant polynomial over $F$, divide by its
leading coefficient to obtain a monic polynomial
$f=T^n+\sum_{i<n}c_iT^i$.  Compute an integer $s\ge0$ such that
$v(c_i)+s(n-i)\ge0$ for every nonzero coefficient.
Then
\[
                    \pi^{sn}f(\pi^{-s}T)\in\mathcal O[T]
                                                               \tag{C.3.GC3}
\]
is monic.  Factor (C.3.GC3) by the preceding procedure and undo
the invertible linear change of variable, normalizing every factor
to be monic.  Restore the original leading coefficient.
Degree-one, nonzero constant and zero polynomials are handled
separately.  This proves the factorization assertion without a
separability restriction on the input polynomial.
\end{proof}

\begin{corollary}[Proved: central division and hermitian arithmetic]
\label{C:Greenberg-valued-algebra-decisions}
Let $F$ satisfy Lemma~\ref{C:Greenberg-congruence-decisions}.
Every explicit finite central simple $F$-algebra has computable
index, actual central division factor, matrix units and an explicit
matrix-algebra isomorphism.  If effective data satisfying that lemma
are supplied uniformly for every valid finite tower field over $F$, their validity
and factorization are decidable, and reduced norms and reduced
characteristic polynomials are computable.

If $\operatorname{char}F\ne2$, let $D$ be an explicit finite-dimensional
division algebra with center $L$ and an explicit $F$-linear involution
$\sigma$ such that $L^\sigma=F$.  For either sign, isometry of
nonsingular hermitian forms over $(D,\sigma)$ is decidable, and
positive decisions construct actual isometries.  The literal
least-Gram and least-pair coordinates, group operations and actual
representatives of Theorem~\ref{C:henselian-hermitian-isometry}
are therefore available under these hypotheses as well.
\end{corollary}
\begin{proof}
We check exactly which parts of the earlier constructions are used.
For a central simple algebra $A$ of dimension $n^2$, the closed
left-ideal scheme $X_j\subseteq\operatorname{Gr}_{nj}(A)$ is given
by (C.3.HP4).  Those homogeneous equations require only finite
linear algebra in the supplied multiplication table.  It is projective,
so Lemma~\ref{C:Greenberg-congruence-decisions} decides each
$X_j(F)\ne\varnothing$.  The elementary left-ideal dimension
calculation in the proof of
Theorem~\ref{C:henselian-central-simple-algorithms} identifies
the least such $j$ with the index $d$.
The same proof then enumerates an idempotent with
$\dim_F Ae=nd$ and a complete set of matrix units.
These are exact finite identities with guaranteed witnesses;
they need neither finite residue nor a completion argument.
The resulting corner $eAe$ and the displayed mutually inverse
matrix maps are the actual asserted division factor and isomorphism.
With the additional uniform tower data, apply the factorization clause
of the lemma successively to decide validity of every tower stage.
The data for a new stage are requested only after its defining
polynomial has passed the preceding stage's irreducibility test.
The finite separable splitting-field search in that earlier proof
then computes reduced determinants and characteristic polynomials.

For hermitian forms of common positive $D$-rank $n$, form
$b=h\perp(-g)$ and the closed projective scheme of right
$D$-stable, totally isotropic subspaces of $D^{2n}$ of half its
$F$-dimension.  Its equations are (C.3.HP4) and (C.3.HP6).
The checked algebraic argument (C.3.HP7), using $2\ne0$ and
hermitian Witt cancellation, says exactly that this scheme has
an $F$-point if and only if $h$ and $g$ are isometric.
Lemma~\ref{C:Greenberg-congruence-decisions} decides that point
question for the actual projective scheme; no smoothness assertion
or residue-field classification is required.
Unequal ranks give a negative answer and common rank zero the
unique zero-space isometry.  For a positive answer the exact
matrix enumeration of the earlier proof constructs an isometry.
Finally, the least-code tests search only the finite set of smaller
codes and use this terminating isometry decision.  Cancellation
turns equality of two Gram pairs into isometry of their cross sums.
Thus the same literal coordinate and operation algorithms terminate.
All finite-residue decisions in the earlier arguments have now been
replaced explicitly by the lemma; the earlier finite-residue
hypothesis is not asserted for the present fields.
\end{proof}
